\subtitle{\bf Aula 16 - Introdução}
\begin{frame}[t,plain]
\titlepage
\end{frame}

\begin{frame}{Objetivos}
\protect\hypertarget{objetivos-1}{}
\begin{itemize}
\item
  Mostrar a correspondência direta entre regras de inferência e equações
  de código Haskell.
\item
  Implementar verificadores de tipos para linguagens com crescente
  complexidade: TExp, TLine, TWhile e TImp.
\end{itemize}
\end{frame}

\hypertarget{introduuxe7uxe3o}{%
\section{Introdução}\label{introduuxe7uxe3o}}

\begin{frame}{Introdução}
\protect\hypertarget{introduuxe7uxe3o-1}{}
\begin{itemize}

\item
  O sistema de tipos define \textbf{o que} é um programa bem tipado via
  regras de inferência
\item
  O \textbf{verificador de tipos} implementa essa especificação como
  código
\end{itemize}
\end{frame}

\begin{frame}{Introdução}
\protect\hypertarget{introduuxe7uxe3o-2}{}
\begin{itemize}

\item
  Correspondência direta: cada regra de inferência vira uma equação
  Haskell
\item
  A estrutura da mônada de verificação espelha os ambientes semânticos
\end{itemize}
\end{frame}

\hypertarget{verificauxe7uxe3o-para-texp}{%
\section{Verificação para TExp}\label{verificauxe7uxe3o-para-texp}}

\begin{frame}[fragile]{A Linguagem TExp}
\protect\hypertarget{a-linguagem-texp}{}
Tipos e termos:

\begin{lstlisting}[language=Haskell]
data Ty   = TBool | TNat

data Term
  = TTrue  | TFalse
  | TIf Term Term Term
  | TZero
  | TSucc Term | TPred Term | TIsZero Term
\end{lstlisting}
\end{frame}

\begin{frame}[fragile]{A Linguagem TExp}
\protect\hypertarget{a-linguagem-texp-1}{}
\begin{itemize}

\item
  Dois tipos: \(\mathbf{Bool}\) e \(\mathbf{Nat}\)
\item
  \textbf{Sem variáveis} -> sem contexto externo necessário
\item
  Mônada: \texttt{TcM = ExceptT String Identity}
\end{itemize}
\end{frame}

\begin{frame}{Da Regra à Equação}
\protect\hypertarget{da-regra-uxe0-equauxe7uxe3o}{}
Regra de inferência:

\[\frac{\vdash t_1 : \mathbf{Bool} \quad \vdash t_2 : T \quad \vdash t_3 : T}{\vdash \mathbf{if}\ t_1\ \mathbf{then}\ t_2\ \mathbf{else}\ t_3 : T} \quad\text{(T-If)}\]
\end{frame}

\begin{frame}[fragile]{Da Regra à Equação}
\protect\hypertarget{da-regra-uxe0-equauxe7uxe3o-1}{}
Equação Haskell:

\begin{lstlisting}[language=Haskell]
tcTerm (TIf t1 t2 t3) = do
    ty1 <- tcTerm t1
    unless (ty1 == TBool) $
        throwError "condition must have type Bool"
    ty2 <- tcTerm t2
    ty3 <- tcTerm t3
    unless (ty2 == ty3) $
        throwError "both branches must have the same type"
    pure ty2
\end{lstlisting}
\end{frame}

\begin{frame}[fragile]{Typechecking}
\protect\hypertarget{typechecking}{}
\begin{lstlisting}[language=Haskell]
tcTerm :: Term -> TcM Ty
tcTerm TTrue  = pure TBool                  -- (T-True)
tcTerm TFalse = pure TBool                  -- (T-False)
tcTerm TZero  = pure TNat                   -- (T-Zero)
\end{lstlisting}
\end{frame}

\begin{frame}[fragile]{Typechecking}
\protect\hypertarget{typechecking-1}{}
\begin{lstlisting}[language=Haskell]
tcTerm (TSucc t1) = do                      -- (T-Succ)
    ty <- tcTerm t1
    unless (ty == TNat) $ throwError "..."
    pure TNat
tcTerm (TIsZero t1) = do                    -- (T-IsZero)
    ty <- tcTerm t1
    unless (ty == TNat) $ throwError "..."
    pure TBool
-- ... TIf, TPred análogos
\end{lstlisting}
\end{frame}

\hypertarget{verificauxe7uxe3o-para-tline}{%
\section{Verificação para TLine}\label{verificauxe7uxe3o-para-tline}}

\begin{frame}[fragile]{A Linguagem TLine}
\protect\hypertarget{a-linguagem-tline}{}
\begin{lstlisting}[language=Haskell]
data Ty = TInt | TBool | TString

data Stmt
  = SDecl   Var Ty Exp   -- var x : T = e ;
  | SAssign Var Exp      -- x := e ;
  | SRead   Exp Var      -- read prompt x ;
  | SPrint  Exp          -- print e ;
\end{lstlisting}
\end{frame}

\begin{frame}[fragile]{Linguagem TLine}
\protect\hypertarget{linguagem-tline}{}
\begin{itemize}

\item
  Principal novidade: \textbf{variáveis} precisam de contexto
\item
  Contexto \(\Gamma : \mathit{Var} \rightharpoonup \mathit{Ty}\)
\item
  Mônada: \texttt{StateT Ctx (ExceptT String Identity)}
\end{itemize}
\end{frame}

\begin{frame}{Linguagem TLine}
\protect\hypertarget{linguagem-tline-1}{}
\begin{itemize}

\item
  Regras para comandos: \(\Gamma \vdash s \dashv \Gamma'\)
\end{itemize}

\[\frac{\Gamma \vdash e : T}{\Gamma \vdash \mathbf{var}\ x : T = e \dashv \Gamma[x \mapsto T]} \quad\text{(T-Decl)}\]
\end{frame}

\begin{frame}{Linguagem TLine}
\protect\hypertarget{linguagem-tline-2}{}
\[\frac{x \in \mathrm{dom}(\Gamma) \quad \Gamma \vdash e : \Gamma(x)}{\Gamma \vdash x \mathbin{:=} e \dashv \Gamma} \quad\text{(T-Assign)}\]

\begin{itemize}

\item
  Declaração \textbf{estende} o contexto:
  \(\Gamma' = \Gamma[x \mapsto T]\)
\item
  Atribuição \textbf{preserva} o contexto: \(\Gamma' = \Gamma\)
\end{itemize}
\end{frame}

\begin{frame}[fragile]{Typechecking}
\protect\hypertarget{typechecking-2}{}
\begin{lstlisting}[language=Haskell]
tcStmt :: Stmt -> TcM ()
tcStmt (SDecl v t e) = do          -- (T-Decl)
    t' <- tcExp e
    unless (t == t') $ throwError "Type error"
    addNewVar v t                   -- Gamma[x -> T]
\end{lstlisting}
\end{frame}

\begin{frame}[fragile]{Typechecking}
\protect\hypertarget{typechecking-3}{}
\begin{lstlisting}[language=Haskell]
tcStmt (SAssign v e) = do
    t  <- lookupVar v
    t' <- tcExp e
    unless (t == t') $ throwError "Type error"
\end{lstlisting}
\end{frame}

\begin{frame}[fragile]{Typechecking}
\protect\hypertarget{typechecking-4}{}
\begin{lstlisting}[language=Haskell]
tcStmt (SRead e v) = do
    t <- tcExp e
    unless (t == TString) $ throwError "Type error"
    t' <- lookupVar v
    unless (t' == TInt) $ throwError "Type error"
tcStmt (SPrint e) = tcExp e >> pure ()
\end{lstlisting}
\end{frame}

\hypertarget{verificauxe7uxe3o-para-twhile}{%
\section{Verificação para TWhile}\label{verificauxe7uxe3o-para-twhile}}

\begin{frame}[fragile]{Escopo}
\protect\hypertarget{escopo}{}
\begin{itemize}

\item
  TWhile adiciona condicionais e repetição com \textbf{escopo local}:
\end{itemize}

\begin{lstlisting}[language=Haskell]
data Stmt
  = -- ... tudo de TLine ...
  | SIf    Exp Block Block
  | SWhile Exp Block

type Block = [Stmt]
\end{lstlisting}
\end{frame}

\begin{frame}[fragile]{Escopo}
\protect\hypertarget{escopo-1}{}
\begin{itemize}

\item
  Variáveis declaradas dentro de
  \texttt{if}/\texttt{while}
  \textbf{não escapam} para fora
\item
  Regras formalizam isso: contexto de saída = contexto de entrada
\end{itemize}
\end{frame}

\begin{frame}{Escopo}
\protect\hypertarget{escopo-2}{}
\[\frac{\Gamma \vdash e : \mathbf{Bool} \quad \Gamma \vdash B_1 \dashv \Gamma_1 \quad \Gamma \vdash B_2 \dashv \Gamma_2}{\Gamma \vdash \mathbf{if}\ e\ \mathbf{then}\ B_1\ \mathbf{else}\ B_2 \dashv \Gamma} \quad\text{(T-If)}\]
\end{frame}

\begin{frame}{Escopo}
\protect\hypertarget{escopo-3}{}
\[\frac{\Gamma \vdash e : \mathbf{Bool} \quad \Gamma \vdash B \dashv \Gamma'}{\Gamma \vdash \mathbf{while}\ e\ \mathbf{do}\ B \dashv \Gamma} \quad\text{(T-While)}\]
\end{frame}

\begin{frame}[fragile]{Implementando Escopo}
\protect\hypertarget{implementando-escopo}{}
\begin{lstlisting}[language=Haskell]
withLocalCtx :: TcM a -> TcM a
withLocalCtx m = do
    ctx <- get
    r   <- m
    put ctx
    pure r
\end{lstlisting}
\end{frame}

\begin{frame}[fragile]{Implementando Escopo}
\protect\hypertarget{implementando-escopo-1}{}
\begin{lstlisting}[language=Haskell]
tcBlock :: [Stmt] -> TcM ()
tcBlock ss = withLocalCtx (mapM_ tcStmt ss)

tcStmt (SIf e b1 b2) = do
    t <- tcExp e
    unless (t == TBool) $ throwError "Type error"
    tcBlock b1
    tcBlock b2
\end{lstlisting}
\end{frame}



\begin{frame}{Ambientes}
\protect\hypertarget{ambientes}{}
\begin{itemize}

\item
  TImp estende TWhile com funções e registros

  \begin{itemize}
  \item
    \begin{description}

    \item[\[\Delta\]]
    ambiente de campos de registros
    \end{description}
  \item
    \begin{description}

    \item[\[\Phi\]]
    ambiente de funções
    \end{description}
  \end{itemize}
\end{itemize}
\end{frame}

\begin{frame}{Relações}
\protect\hypertarget{relauxe7uxf5es}{}
\begin{itemize}

\item
  Expressões: \(\Gamma;\Delta;\Phi \vdash e : \tau\)
\item
  Comandos:
  \(\Gamma;\Delta;\Phi;\rho_\mathit{ret} \vdash s \dashv \Gamma'\)
\end{itemize}
\end{frame}

\begin{frame}{Registros}
\protect\hypertarget{registros}{}
\begin{itemize}

\item
  \textbf{Acesso a campo (T-Field):}
\end{itemize}

\[\frac{\Gamma;\Delta;\Phi \vdash e : R \quad \Delta(R)(f) = \tau}{\Gamma;\Delta;\Phi \vdash e.f : \tau}\]
\end{frame}

\begin{frame}{Registros}
\protect\hypertarget{registros-1}{}
\begin{itemize}

\item
  \textbf{Construção de registro (T-New):}
\end{itemize}

\[\frac{R \in \mathrm{dom}(\Delta) \quad \Delta(R) = \overline{(f_i:\tau_i)} \quad \Gamma;\Delta;\Phi \vdash e_i : \tau_i}{\Gamma;\Delta;\Phi \vdash \mathbf{new}\ R\{\overline{f_i=e_i}\} : R}\]
\end{frame}

\begin{frame}{Funções}
\protect\hypertarget{funuxe7uxf5es}{}
\begin{itemize}

\item
  \textbf{Chamada como expressão (T-CallExpr):}
\end{itemize}

\[\frac{\Phi(f) = ([\tau_1,\ldots,\tau_n],\;\tau) \quad \Gamma;\Delta;\Phi \vdash e_i : \tau_i}{\Gamma;\Delta;\Phi \vdash f(e_1,\ldots,e_n) : \tau}\]
\end{frame}

\begin{frame}{Funções}
\protect\hypertarget{funuxe7uxf5es-1}{}
\begin{itemize}

\item
  \textbf{Declaração de função (T-FuncDecl):}
\end{itemize}

\[\frac{[x_1\!\mapsto\!\tau_1,\ldots,x_n\!\mapsto\!\tau_n];\Delta;\Phi;\rho \vdash B \dashv \_}{\Delta;\Phi \vdash \mathbf{fn}\ f(\overline{x_i:\tau_i}):\rho\ \{B\}\ \checkmark}\]
\end{frame}

\hypertarget{verificauxe7uxe3o-para-timp}{%
\section{Verificação para TImp}\label{verificauxe7uxe3o-para-timp}}

\begin{frame}[fragile]{A Mônada}
\protect\hypertarget{a-muxf4nada}{}
\begin{itemize}

\item
  TImp adiciona funções e registros --- Informação \textbf{imutável} vai
  em \texttt{ReaderT}:
\end{itemize}

\begin{lstlisting}[language=Haskell]
data TcEnv = TcEnv
  { recordEnv  :: Map Name [(Field, Ty)]  -- Delta
  , funcEnv    :: Map Name ([Ty], RetTy)  -- Phi
  , returnType :: Maybe RetTy             -- rho_ret
  }
\end{lstlisting}
\end{frame}

\begin{frame}[fragile]{A Mônada}
\protect\hypertarget{a-muxf4nada-1}{}
\begin{lstlisting}[language=Haskell]
type VarCtx = Map Var Ty

type TcM a = ReaderT TcEnv
               (StateT VarCtx
                 (ExceptT String Identity)) a
\end{lstlisting}
\end{frame}

\begin{frame}[fragile]{Registros}
\protect\hypertarget{registros-2}{}
\begin{lstlisting}[language=Haskell]
tcExp (EField e f) = do
  t <- tcExp e
  case t of
    TRecord rname -> lookupField rname f
    _ -> throwError "Field access on non-record type"
\end{lstlisting}
\end{frame}

\begin{frame}[fragile]{Funções}
\protect\hypertarget{funuxe7uxf5es-2}{}
\begin{lstlisting}[language=Haskell]
tcExp (ECall fname args) = do
  fenv <- asks funcEnv
  case Map.lookup fname fenv of
    Nothing -> throwError ("Undefined function: " ++ fname)
    Just (pts, rt) -> do
      argTys <- mapM tcExp args
      unless (argTys == pts) $ throwError "Argument type mismatch"
      case rt of
        RTVoid  -> throwError "Void function used as expression"
        RTTy ty -> pure ty
\end{lstlisting}
\end{frame}

\begin{frame}[fragile]{Função}
\protect\hypertarget{funuxe7uxe3o}{}
\begin{lstlisting}[language=Haskell]
tcFuncDecl :: TcEnv -> FuncDecl -> Either String ()
tcFuncDecl baseEnv (FuncDecl name params retTy body) = do
  let paramPairs = [(v, t) | Param v t <- params]
      initCtx    = Map.fromList paramPairs
      env        = baseEnv { returnType = Just retTy }
  case runIdentity (runExceptT
         (execStateT (runReaderT (tcBlock body) env) initCtx)) of
    Left err -> Left ("In function " ++ name ++ ": " ++ err)
    Right _  -> Right ()
\end{lstlisting}
\end{frame}

\begin{frame}[fragile]{Função}
\protect\hypertarget{funuxe7uxe3o-1}{}
\begin{itemize}

\item
  Contexto inicial = parâmetros formais
\item
  \texttt{returnType} fixado via
  \texttt{local} para verificar
  \texttt{return}
\item
  Variáveis externas à função são \textbf{invisíveis}
\end{itemize}
\end{frame}

\hypertarget{conclusuxe3o}{%
\section{Conclusão}\label{conclusuxe3o}}

\begin{frame}{Conclusão}
\protect\hypertarget{conclusuxe3o-1}{}
\begin{itemize}

\item
  O verificador de tipos é uma \textbf{transcrição sistemática} do
  sistema de tipos
\item
  Cada regra de inferência corresponde a uma equação Haskell
\end{itemize}
\end{frame}
